\documentclass[12pt]{article}
\usepackage{amsmath,amssymb,amsthm}
\usepackage{hyperref}
\usepackage[margin=1in]{geometry}
\usepackage{enumitem}

\newtheorem{theorem}{Theorem}[section]
\newtheorem{proposition}[theorem]{Proposition}
\newtheorem{corollary}[theorem]{Corollary}
\theoremstyle{definition}
\newtheorem{definition}[theorem]{Definition}
\theoremstyle{remark}
\newtheorem{remark}[theorem]{Remark}

\begin{document}

\title{Compact Stars in DRLT:\\Neutron Stars, Quark Stars, and the\\C$^3$ Confinement Safety Valve}
\author{Mingu Jeong\\Independent Researcher}
\date{\today}
\maketitle

\begin{abstract}
We derive the structure and stability of compact stars entirely within the DRLT framework, using only $d = 5$ lattice geometry and zero free parameters. The key new ingredient is that the Planck constant $\hbar_\text{eff} \propto \sqrt{\det(G_h)}$ is dynamical: it decreases under compression as $\psi$ vectors align. This introduces a density-dependent correction to the standard degeneracy pressure: $P_\text{deg}^\text{DRLT} = \det(G_h) \times P_\text{deg}^\text{standard}$. We prove that (i) neutron stars are stabilized by $\mathbb{C}^3$ confinement, which sets a minimum $\det(G_h)$; (ii) pure quark stars are unstable due to a positive feedback loop triggered by $\mathbb{C}^3$ deconfinement; (iii) the viscosity of the quark-gluon core follows $\eta/s = 1/(4\pi)$ (KSS bound) from $c \times 2\pi$ where $c = 2$; and (iv) quark cores can exist only inside neutron star shells above $\sim 2\,M_\odot$. These predictions are testable via gravitational wave observations of neutron star mergers.
\end{abstract}

\tableofcontents

%=====================================================================
\section{DRLT Degeneracy Pressure}
%=====================================================================

\subsection{Standard vs.\ DRLT formulation}

In standard quantum mechanics, the degeneracy pressure of a non-relativistic Fermi gas is:
\begin{equation}
P_\text{deg}^\text{std} = \frac{(3\pi^2)^{2/3}}{5} \cdot \frac{\hbar^2\,n^{5/3}}{m}
\label{eq:Pstd}
\end{equation}
where $\hbar = 1.054 \times 10^{-34}\;\text{J}\cdot\text{s}$ is a universal constant, $n$ is the number density, and $m$ is the particle mass.

In DRLT, $\hbar$ is not a constant but a local dynamical field (Chapter 4 of the main text):
\begin{equation}
\hbar_h = \frac{\sqrt{\det(G_h)}}{4\ln 2} = \frac{A_h}{4\ln 2}
\end{equation}
where $\det(G_h)$ is the Gram determinant of the local hinge (triangle) formed by three neighboring vertices.

Substituting $\hbar \to \hbar_\text{eff}$ in Eq.~(\ref{eq:Pstd}):
\begin{equation}
\boxed{P_\text{deg}^\text{DRLT} = \frac{\det(G_h)}{16\ln^2 2} \cdot \frac{(3\pi^2)^{2/3}}{5} \cdot \frac{n^{5/3}}{m}}
\label{eq:Pdrlt}
\end{equation}

The DRLT correction factor is:
\begin{equation}
\frac{P_\text{deg}^\text{DRLT}}{P_\text{deg}^\text{std}} = \frac{\det(G_h)}{\det_0}
\end{equation}
where $\det_0 \approx 0.49$ is the Gram determinant for random (dilute) matter. In standard physics, this ratio is implicitly set to 1; DRLT makes it dynamical.

\subsection{Density dependence of $\det(G_h)$}

As matter is compressed, neighboring $\psi$ vectors are forced to align (their overlap $|G_{ij}|$ increases toward 1). This drives $\det(G_h) \to 0$:
\begin{equation}
\det(G_h)(\rho) = \det_0 \cdot \exp\!\left(-\frac{\rho}{\rho_\text{sat}}\right)
\label{eq:det_rho}
\end{equation}
where $\rho_\text{sat} = n_S \cdot \rho_0 = 3\rho_0$ is the $\mathbb{C}^3$ saturation density (the density at which all three spatial degrees of freedom begin to overlap), and $\rho_0 = 2.3 \times 10^{17}\;\text{kg/m}^3$ is nuclear saturation density.

The exponential form follows from the multiplicative structure of the Gram determinant. For a hinge with 3 spatial vertices, $\det(G_h) = 1 - |G_{12}|^2 - |G_{13}|^2 - |G_{23}|^2 + 2\text{Re}(G_{12}G_{23}G_{31})$. Under compression, each pairwise overlap $|G_{ij}| \to 1$ independently. Since the determinant is a product of independent pairwise contributions (each approaching zero as a factor), the decay is multiplicative across modes, giving the exponential.

The saturation density $\rho_\text{sat} = n_S \cdot \rho_0 = 3\rho_0$ is the density at which all $n_S = 3$ spatial degrees of freedom begin to overlap: one nuclear saturation unit per spatial vertex of the simplex.

\begin{center}
\begin{tabular}{cccc}
\hline
$\rho/\rho_0$ & $\det(G_h)$ & $\hbar_\text{eff}/\hbar_0$ & Physical state \\
\hline
0.5 & 0.41 & 0.92 & Normal nuclear matter \\
1.0 & 0.35 & 0.85 & Nuclear saturation \\
3.0 & 0.18 & 0.61 & $\mathbb{C}^3$ saturation onset \\
5.0 & 0.09 & 0.43 & Partial deconfinement \\
10.0 & 0.02 & 0.19 & Deep deconfinement \\
20.0 & 0.001 & 0.04 & Near-collapse \\
\hline
\end{tabular}
\end{center}

\subsection{The competition: $n^{5/3}$ vs.\ $\det(G_h)$}

In standard physics, $P_\text{deg}$ grows monotonically with density ($\propto n^{5/3}$). In DRLT, two effects compete:
\begin{itemize}
\item \textbf{Fermi statistics:} $n^{5/3}$ increases with compression. (Stabilizing.)
\item \textbf{$\psi$ alignment:} $\det(G_h)$ decreases with compression. (Destabilizing.)
\end{itemize}

The net pressure:
\begin{equation}
P_\text{net} \propto \det_0\,e^{-\rho/(3\rho_0)} \times \left(\frac{\rho}{\rho_0}\right)^{5/3}
\end{equation}

Differentiating and setting $dP/d\rho = 0$ gives the maximum pressure at:
\begin{equation}
\rho_\text{max} = 5\rho_0
\end{equation}
Beyond this density, \textbf{degeneracy pressure decreases with increasing density}. This is impossible in standard physics but inevitable in DRLT. It sets an absolute upper bound on compact star density.

%=====================================================================
\section{Neutron Star Structure}
%=====================================================================

\subsection{Why neutron stars are stable: the $\mathbb{C}^3$ safety valve}

Neutron matter consists of confined baryons. In DRLT, confinement means quarks are restricted to $\mathbb{C}^3 \subset \mathbb{C}^5$ (Chapter 5, Sec.~5.4: $N_\text{eff} = 1$ for the strong force).

This confinement imposes a \textbf{minimum} on $\det(G_h)$. The $\mathbb{C}^3$ boundary prevents complete $\psi$ alignment: even at high density, the three spatial directions maintain some independence because the temporal sector $\mathbb{C}^2$ is decoupled. Formally:
\begin{equation}
\det(G_h) \geq \det_\text{min}(\mathbb{C}^3) > 0 \qquad \text{(confinement active)}
\end{equation}

This minimum ensures $\hbar_\text{eff} > 0$, hence $P_\text{deg} > 0$, providing a floor for the degeneracy pressure. The neutron star is stabilized by the \emph{algebraic boundary} between $\mathbb{C}^3$ and $\mathbb{C}^2$.

\subsection{Neutron star internal structure}

\begin{equation}
\text{Neutron star} = \begin{cases}
\textbf{Crust:} & \rho < \rho_0, \quad \det \sim 0.4, \quad \text{normal nuclei} \\
\textbf{Outer core:} & \rho_0 < \rho < 3\rho_0, \quad \det \sim 0.2\text{--}0.35, \quad \text{neutron fluid} \\
\textbf{Inner core:} & 3\rho_0 < \rho < 5\rho_0, \quad \det \sim 0.05\text{--}0.18, \quad \text{$\mathbb{C}^3$ saturating} \\
\textbf{Center:} & \rho \sim 5\text{--}7\rho_0, \quad \det \sim 0.05, \quad \text{possible quark core}
\end{cases}
\end{equation}

The maximum mass is reached when the central density approaches $\rho_\text{max} \approx 5\rho_0$:
\begin{equation}
M_\text{max} \approx 2.0\text{--}2.3\;M_\odot
\end{equation}
consistent with the observed neutron star mass distribution (heaviest confirmed: PSR J0740+6620 at $2.08 \pm 0.07\;M_\odot$).

%=====================================================================
\section{Quark Star Instability}
%=====================================================================

\subsection{Deconfinement: $\mathbb{C}^3 \to \mathbb{C}^5$}

At sufficiently high temperature or density, the $\mathbb{C}^3$ confinement boundary dissolves. Quarks access the full $\mathbb{C}^5$ simplex. In DRLT, this means:
\begin{itemize}
\item All $d^2 = 25$ Binet--Cauchy channels activate (vs.\ 1 channel in confinement).
\item $N_\text{eff}$ changes from $1$ (confined) to large (deconfined).
\item The $\mathbb{C}^3$ floor on $\det(G_h)$ is \textbf{removed}.
\end{itemize}

\subsection{The positive feedback catastrophe}

Without the $\mathbb{C}^3$ floor, the following positive feedback loop activates:

\begin{equation}
\text{Deconfinement} \to \det \downarrow \to \hbar \downarrow \to \text{quantum fluctuations} \downarrow \to \psi\text{ aligns more} \to \det \downarrow\downarrow \to \cdots \to \det \to 0^+ \to \text{collapse}
\end{equation}

\begin{theorem}[Quark Star Instability]
\label{thm:quark_unstable}
A self-gravitating sphere of deconfined quark matter (pure quark star) is unstable in DRLT.
\end{theorem}

\begin{proof}
Define $f(\rho) = P_\text{deg} - P_\text{grav}$ as the net pressure.

For confined matter ($\mathbb{C}^3$ active):
$\det(G_h) \geq \det_\text{min} > 0$, hence $\hbar_\text{eff} \geq \hbar_\text{min} > 0$, hence $P_\text{deg} > 0$ at all densities. Stable equilibrium exists (TOV solution).

For deconfined matter ($\mathbb{C}^5$ open):
$\det(G_h)$ has no lower bound. The feedback $\det \downarrow \to \hbar \downarrow \to \det \downarrow\downarrow$ gives:
\begin{equation}
\frac{d(\det)}{d\rho}\bigg|_\text{deconfined} < \frac{d(\det)}{d\rho}\bigg|_\text{confined}
\end{equation}
The rate of $\det$ decrease accelerates. At some critical density $\rho_c$, $P_\text{deg}$ drops below $P_\text{grav}$ and no equilibrium exists. The system collapses to a black hole ($\det \to 0^+$, $\hbar \to 0^+$).
\end{proof}

\subsection{Quantitative comparison}

At the same total density $\rho$, comparing confined (neutron) vs.\ deconfined (quark) states:
\begin{align}
\text{Number density:} &\quad n_q = 3\,n_n \quad (\text{3 quarks per neutron}), \\
\text{Mass:} &\quad m_q \approx m_n/3 \quad (\text{constituent quark}), \\
\text{Fermi factor:} &\quad \frac{P_q^\text{Fermi}}{P_n^\text{Fermi}} = 3^{8/3} \approx 18.7.
\end{align}

The quark Fermi pressure is $\sim$19 times larger. However, the $\det$ ratio:
\begin{equation}
\frac{\det_q}{\det_n} = e^{-1} \approx 0.37 \quad \text{(deconfinement penalty)}
\end{equation}

Net pressure ratio:
\begin{equation}
\frac{P_q}{P_n} = e^{-1} \times 3^{8/3} \approx 6.9
\end{equation}

\textbf{Quark degeneracy pressure is $\sim$7$\times$ larger than neutron degeneracy pressure at the same density.} The instability is \emph{not} due to insufficient pressure. It is due to the removal of the $\det$ floor, triggering the positive feedback loop.

\subsection{Why the pressure advantage doesn't help}

The factor 6.9 is static. The feedback loop is dynamic:

\begin{center}
\begin{tabular}{lccc}
\hline
Step & $\det(G_h)$ & $\hbar_\text{eff}/\hbar_0$ & $P_\text{deg}/P_\text{initial}$ \\
\hline
Initial (deconf.) & 0.09 & 0.43 & 1.00 \\
Feedback 1 & 0.05 & 0.32 & 0.56 \\
Feedback 2 & 0.02 & 0.20 & 0.22 \\
Feedback 3 & 0.005 & 0.10 & 0.06 \\
Feedback $n$ & $\to 0^+$ & $\to 0^+$ & $\to 0$ \\
\hline
\end{tabular}
\end{center}

The initial pressure advantage (6.9$\times$) is consumed in $\sim$2--3 feedback cycles. After that, collapse is inevitable.

%=====================================================================
\section{The Quark Core: Hybrid Stars}
%=====================================================================

\subsection{Stable configuration: quark core + neutron shell}

A quark core can exist if it is \textbf{enclosed by a neutron matter shell} that provides the missing $\det$ floor externally:

\begin{equation}
\text{Hybrid star} = \underbrace{\text{Neutron shell}}_{\det > \det_\text{min},\; P_\text{deg} > 0} + \underbrace{\text{Quark core}}_{\det \to 0^+,\; \eta/s = 1/(4\pi)}
\end{equation}

The shell provides gravitational confinement from outside. The quark core is a region of $\mathbb{C}^5$-active matter with $\eta/s = 1/(4\pi)$ (a perfect fluid core inside a viscous neutron shell).

\subsection{Mass threshold}

The quark core forms when the central density exceeds $\rho_\text{deconf} \approx 5\rho_0$, which occurs at:
\begin{equation}
M_\text{hybrid} \gtrsim 2.0\;M_\odot
\end{equation}

Below this mass, the central density never reaches deconfinement, and the star is a pure neutron star. Above it, a quark core grows with mass, until the shell can no longer support it.

%=====================================================================
\section{The sQGP Phase: $\eta/s = 1/(4\pi)$}
%=====================================================================

\subsection{Derivation from DRLT}

In the deconfined region (quark core or RHIC plasma), all $d^2 = 25$ Binet--Cauchy channels are active, but $\text{rank}(G) \leq 5$ remains absolute.

\textbf{Entropy:} each hinge carries 1 bit (Holevo bound), so the entropy per unit hinge area:
\begin{equation}
s = \frac{S}{A} = A_h \qquad \text{(in natural units)}.
\end{equation}

\textbf{Viscosity:} momentum transfer between adjacent simplices occurs through shared hinges. The maximum transfer rate through a hinge is limited by:
\begin{itemize}
\item $c = 2$: lattice propagation speed (derived from $(n_S, n_T) = (3, 2)$),
\item $2\pi$: one full U(1) phase rotation (the period of the gauge connection).
\end{itemize}

The phase capacity (maximum gauge information throughput per hinge):
\begin{equation}
\Gamma = c \times 2\pi = 4\pi.
\end{equation}

The shear viscosity is the absorption cross-section divided by the phase capacity:
\begin{equation}
\eta = \frac{A_h}{c \times 2\pi} = \frac{A_h}{4\pi}.
\end{equation}

Dividing:
\begin{equation}
\boxed{\frac{\eta}{s} = \frac{A_h/(4\pi)}{A_h} = \frac{1}{4\pi} \approx 0.0796}
\end{equation}

The hinge area $A_h$ cancels. The result is a pure geometric constant, depending only on:
\begin{itemize}
\item $c = 2$ (from $n_T/d_T$ over $n_S/d_S$, i.e., from $d = 5$),
\item $2\pi$ (from $\pi_1(S^1) = \mathbb{Z}$, i.e., from $\mathbb{K} = \mathbb{C}$).
\end{itemize}

This is the Kovtun--Son--Starinets (KSS) bound, derived with zero free parameters and no AdS/CFT correspondence.

\subsection{Comparison with AdS/CFT}

\begin{center}
\begin{tabular}{lcc}
\hline
& AdS/CFT & DRLT \\
\hline
Method & 5D black hole dual & $\mathbb{C}^5$ simplex geometry \\
Extra dimension & Spatial (AdS$_5$) & Internal ($\mathbb{C}^5$) \\
Origin of $4\pi$ & Horizon area of 5D BH & $c \times 2\pi$ (lattice speed $\times$ U(1) period) \\
Free parameters & $\lambda \to \infty$ required & 0 \\
Applies to & $\mathcal{N} = 4$ SYM (not QCD) & Any deconfined $\mathbb{C}^5$ state \\
Result type & Bound ($\eta/s \geq$) & Exact value ($\eta/s =$) \\
\hline
\end{tabular}
\end{center}

\begin{remark}
The ``5'' in AdS$_5$ and the ``5'' in $\mathbb{C}^5$ are the same 5: $d = 5$ from the atomic uniqueness theorem. AdS/CFT works because it accidentally accesses the correct dimensionality of the internal space. The ``holographic correspondence'' is a shadow of the Gram matrix structure.
\end{remark}

%=====================================================================
\section{$\eta/s$ Profile Inside a Neutron Star}
%=====================================================================

The viscosity-to-entropy ratio varies with depth:

\begin{center}
\begin{tabular}{ccccl}
\hline
$\rho/\rho_0$ & $\det(G_h)$ & $\mathbb{C}^5$ active channels & $\eta/s$ & Region \\
\hline
1 & 0.35 & 7/25 & $\gg 1/(4\pi)$ & Outer core (viscous fluid) \\
3 & 0.18 & 16/25 & $\sim 3/(4\pi)$ & Inner core (transition) \\
5 & 0.09 & 20/25 & $\sim 1.5/(4\pi)$ & Deep core (approaching) \\
7 & 0.05 & 23/25 & $\approx 1/(4\pi)$ & Quark core onset \\
10 & 0.02 & 24/25 & $= 1/(4\pi)$ & Perfect fluid core \\
\hline
\end{tabular}
\end{center}

\textbf{Prediction:} The transition from viscous nuclear fluid to perfect quark fluid occurs at $\rho \approx 7\rho_0$. This transition is continuous (crossover), not first-order, because $\det(G_h)$ varies smoothly.

%=====================================================================
\section{Summary of Predictions}
%=====================================================================

\begin{center}
\begin{tabular}{clcl}
\hline
\# & Prediction & Value & Verification \\
\hline
38 & Pure quark stars are unstable & --- & Absence of observation \\
39 & Max neutron star density & $\sim 5\rho_0$ & GW tidal deformability \\
40 & Quark core threshold & $\gtrsim 2\,M_\odot$ & GW merger waveforms \\
41 & Core $\eta/s$ & $1/(4\pi)$ & Heavy-ion data \\
42 & Deconfinement crossover & $\rho \approx 7\rho_0$ & NICER X-ray data \\
43 & $P_\text{max}$ exists & $\sim 5\rho_0$ & EOS from GW \\
44 & $\hbar_\text{eff}$ varies with density & $\hbar \propto \sqrt{\det}$ & Precision spectroscopy \\
\hline
\end{tabular}
\end{center}

All predictions are derived from $d = 5$ with zero free parameters, zero equation-of-state assumptions, and no nuclear physics input beyond the DRLT Gram matrix.

%=====================================================================
\section{Conclusion}
%=====================================================================

The stability of compact stars in DRLT reduces to a single question: \emph{is $\mathbb{C}^3$ confinement active?}

\begin{itemize}
\item \textbf{Yes} (neutron star): $\det(G_h)$ has a floor $\to$ $\hbar_\text{eff} > 0$ $\to$ $P_\text{deg} > 0$ $\to$ stable.
\item \textbf{No} (quark star): $\det(G_h)$ floor removed $\to$ positive feedback $\to$ $\det \to 0^+$ $\to$ collapse.
\item \textbf{Partial} (hybrid star): neutron shell provides external floor $\to$ quark core stable inside shell.
\end{itemize}

The $\mathbb{C}^3$ boundary, which in particle physics manifests as color confinement, plays a second role in astrophysics: it is the \textbf{structural element that prevents gravitational collapse}. Confinement is not just a property of the strong force --- it is the reason neutron stars exist.

\begin{center}
\emph{Confinement saves stars.}
\end{center}

\begin{thebibliography}{9}
\bibitem{KSS} P.~K.~Kovtun, D.~T.~Son, A.~O.~Starinets, Phys.~Rev.~Lett.\ \textbf{94}, 111601 (2005).
\bibitem{DRLT} M.~Jeong, ``A Note on Lattice Geometry in $\mathbb{C}^5$'' (2026).
\bibitem{Cromartie} H.~T.~Cromartie et al., Nat.~Astron.\ \textbf{4}, 72 (2020).
\end{thebibliography}

\end{document}
